% !TeX root = ../../Book3.tex
% 纲要标题：## 第18章　正则局部环
% 纲要位置：计划纲要.md，第 2519 行。

\chapter{正则局部环}
\label{b3:ch:18}
\BookChapterNumber{b3:ch:18}

\begin{chapterabstract}
本章建立正则局部环的数值、伴随分次和同调刻画，并由此研究局部化、完备化及唯一分解性。剩余域的有限长度、各项有限秩的自由消解与极大理想的正则参数系给出两种不同检验；环论正则性、相对于基域的光滑性与几何正则性也得到区分。
\end{chapterabstract}

\begin{learninggoals}
\begin{itemize}
\item 在极大理想生成元、伴随分次环和正则序列之间转换正则性条件。
\item 证明剩余域的有限投射维数判据，并把有限生成模的投射维数上界推广到任意模。
\item 运用消解局部化和完备化的有限商比较，判断正则性的保持。
\item 说明有限自由消解怎样使主开集上的可逆模自由，并完成唯一分解证明。
\item 通过切空间和不可分基变换的计算，区分正则性、相对光滑性与几何正则性。
\end{itemize}
\end{learninggoals}

设 \(k\) 为域。尖点环 \(k[t^2,t^3]\) 是一维整环；在原点对应的局部环中，极大理想由 \(t^2,t^3\) 两个元素生成，不能缩成一个。相比之下，直线原点的局部环只有一个局部参数。两者的素理想链长度都是一，差别却出现在需要多少个独立的一阶生成元。正则局部环要求这两个数相等，因此它比“没有零因子”强得多。

\ProseParagraphBreak

这个数值等式若只是定义，还不能解释为什么正则局部环如此接近多项式环。应当进一步考察极大理想逐层的商：若各层像自由变量的齐次多项式一样增长，参数在逐次商中便没有意外关系。Koszul 复形和剩余域的自由消解又从同调方向检验同一现象。把这些判据连接起来，才能证明局部化、完备化与唯一分解的后续性质。

\ProseParagraphBreak

本章需要\cref{b3:thm:hilbert-samuel}、\cref{b3:thm:krull-intersection}、\cref{b3:thm:koszul-acyclic}和\cref{b3:thm:auslander-buchsbaum}。整体维数部分使用第二卷第7.2节已经证明的 Baer 判据\footnote{贝尔（Reinhold Baer，1902-1979），德国数学家，研究群论并引入内射模的概念。}与足够多内射模定理；唯一分解性部分使用第二卷第9.2—9.3节的外幂与行列式。光滑性的讨论使用\chapref{b3:ch:14}的定义、Jacobian 修正及不可分例子。

\ProseParagraphBreak
同调判据及正则性的局部化可参阅松村英之的《Commutative Ring Theory》\cite[Theorems 19.2--19.3, pp. 156--157]{Matsumura1989CommutativeRingTheory}。高度一素理想判据和唯一分解性证明见同书\cite[Theorems 20.1--20.3, pp. 161--164]{Matsumura1989CommutativeRingTheory}；本章把秩一模的有限展示延拓和行列式步骤分别展开。

% !TeX root = ../../Book3.tex
% 纲要标题：### 18.1 正则性的等价观点
% 纲要位置：工作记录/计划纲要.md，第 2552 行。
% 纲要细目（写作范围）：
% * 维数与极大理想最小生成元数
% * 以规范映射 \(\operatorname{Sym}_k(\mathfrak m/\mathfrak m^2)\to\operatorname{gr}_{\mathfrak m}R\) 的同构性独立刻画正则性；反向比较 Hilbert–Samuel 增长次数
% * 逐层说明正则参数商的伴随分次同构；证明任意参数系正则时写出商环中的根等式
% * 正文命题完整计算图形环、二次锥及 \(\mathbb Z[x]_{(p,x)}\) 的正则性、伴随分次环和混合特征剩余域消解

\section{正则性的等价观点}
\label{b3:sec:1801}

极大理想的生成元数等于维数，是正则局部环的一个数值条件。伴随分次环把这个等式转为具体的多项式结构，由此可以证明极小生成元构成正则序列。

% 纲要内容（仅作写作计划，尚非教材正文）：
%
% * 维数与极大理想最小生成元数
% * 以规范映射 \(\operatorname{Sym}_k(\mathfrak m/\mathfrak m^2)\to\operatorname{gr}_{\mathfrak m}R\) 的同构性独立刻画正则性；反向比较 Hilbert–Samuel 增长次数
% * 逐层说明正则参数商的伴随分次同构；证明任意参数系正则时写出商环中的根等式
%

\subsection{维数与极大理想的最小生成元数}
\label{b3:subsec:1801:01}

\cref{b3:def:regular-local-ring}以 \(\dim R=\operatorname{edim}R\) 定义正则局部环。这个等式说：极大理想所需的生成元恰好达到高度定理允许的最少数目。本章要从这一个数字等式推出伴随分次环的多项式结构、有限自由消解以及唯一分解。

\ProseParagraphBreak
设 \((R,\mathfrak m)\) 为交换 Noether 局部环，\(k=R/\mathfrak m\)，取 \(\mathfrak m\) 的极小生成组 \(x_1,\ldots,x_e\)。它们在 \(\mathfrak m/\mathfrak m^2\) 中的像为基，从而给出分次满射
\begin{equation}\label{b3:eq:regular-graded-presentation}
k[X_1,\ldots,X_e]\longrightarrow\operatorname{gr}_{\mathfrak m}R,
\qquad X_i\longmapsto x_i+\mathfrak m^2.
\end{equation}
因为 \(\mathfrak m=(x_1,\ldots,x_e)\)，\(\mathfrak m^n\) 中的元素都是 \(n\) 个生成元之积的 \(R\) 线性组合；模掉 \(\mathfrak m^{n+1}\) 后，系数只需保留它们在 \(k\) 中的像。因此每一层都由 \(x_i+\mathfrak m^2\) 的乘积生成，得到上述满射。满射性并不依赖正则性。不选基时，对称代数的泛性质仍把一次分量的恒等映射延拓为规范同态 \(\operatorname{Sym}_k(\mathfrak m/\mathfrak m^2)\to\operatorname{gr}_{\mathfrak m}R\)，选基以后才写成\eqref{b3:eq:regular-graded-presentation}。一般局部环可能存在初始关系；正则性将使这个核恰好为零。

\begin{theorem}\label{b3:thm:regular-local-characterizations}
设 \((R,\mathfrak m)\) 为交换 Noether 局部环，\(k=R/\mathfrak m\)、\(d=\dim R\)，并记 \(\operatorname{gr}_{\mathfrak m}R=\bigoplus_{q\ge0}\mathfrak m^q/\mathfrak m^{q+1}\)。下列条件等价：
\begin{equivalences}
\item \(R\) 正则，即 \(\dim_k\mathfrak m/\mathfrak m^2=d\)；
\item 由一次分量 \(\mathfrak m/\mathfrak m^2\) 的恒等映射诱导的规范分次 \(k\) 代数同态
\[
\operatorname{Sym}_k(\mathfrak m/\mathfrak m^2)
\longrightarrow\operatorname{gr}_{\mathfrak m}R
\]
为同构；
\item \(\mathfrak m\) 可由一个 \(R\) 正则序列生成。
\end{equivalences}
这些条件成立时，任意极小生成组都是长度 \(d\) 的正则序列，\(\operatorname{depth}R=d\)；任意参数系也是正则序列。
\end{theorem}

参数系只要求生成极大理想准素的理想，未必生成极大理想本身。极小生成组与一般参数系的正则性因而需要分别证明。

\subsection{伴随分次环的结构}
\label{b3:subsec:1801:02}

\begin{proof}[\cref{b3:thm:regular-local-characterizations}的证明]
假设 \(R\) 正则，在\eqref{b3:eq:regular-graded-presentation}中取 \(e=d\)。\(d=0\) 时 Nakayama 引理给 \(\mathfrak m=0\)，满射本来就是域的恒等。以下设 \(d\ge1\)。满射是否单射，可以从各层的增长判断：\(d\) 个变量的多项式环具有 \(d\) 次累计增长，而一个非零初始关系就会使这个次数降低。若核含非零齐次多项式 \(f\)，记其次数为 \(a>0\)；次数零的映射是 \(k\) 的恒等，故非零关系不可能在次数零。多项式环中的乘 \(f\) 单射，把次数至多 \(n-a\) 的部分送到次数至多 \(n\) 的部分，故其模 \(f\) 商环的前 \(n\) 层累计维数为
\[
\binom{n+d}{d}-\binom{n-a+d}{d}\qquad(n\ge a),
\]
这两个二项式多项式的最高项都是 \(n^d/d!\)，相减后抵消，故结果的次数至多为 \(d-1\)。伴随分次环是该商环的商，其累计维数不超过这个数。但累计维数等于 \(\ell_R(R/\mathfrak m^{n+1})\)，由\cref{b3:thm:hilbert-samuel}恰有 \(d\) 次正首项，不可能在充分大的 \(n\) 处被一个至多 \(d-1\) 次的多项式控制，矛盾。因此满射总为同构。选定这组基后，对称代数识别为 \(k[X_1,\ldots,X_d]\)，上述满射就是条件中的规范同态。

反向，设该规范同态为同构，并记 \(e=\dim_k\mathfrak m/\mathfrak m^2\)。选取一次分量的一组基，就把伴随分次环识别为 \(k[X_1,\ldots,X_e]\)，因而对每个 \(n\ge0\)，
\[
\ell_R(R/\mathfrak m^{n+1})
=\sum_{q=0}^n\dim_k\mathfrak m^q/\mathfrak m^{q+1}
=\binom{n+e}{e}.
\]
右端是 \(e\) 次多项式；Hilbert–Samuel 定理给出的次数为 \(d\)，所以 \(e=d\)，即 \(R\) 正则。

现在假设规范分次同态为同构。已知 \(e=d\)，取任意极小生成组 \(x_1,\ldots,x_d\)，其对应的分次同构使 \(x_i\) 的初始式为 \(X_i\)。\(d=0\) 时 \(\mathfrak m=0\)，由空正则序列生成；以下设 \(d\ge1\)。要把分次环中的非零因子转回 \(R\)，须保证每个非零元素都能被某一层检测到。由\cref{b3:thm:krull-intersection}，每个非零 \(a\in R\) 都有最大的整数 \(\nu(a)\ge0\) 使 \(a\in\mathfrak m^{\nu(a)}\)，其初始式非零。\(x_1a\) 在 \(\mathfrak m^{\nu(a)+1}/\mathfrak m^{\nu(a)+2}\) 中的像为 \(X_1\operatorname{in}(a)\)，在多项式环中非零；所以乘积的最低项不会消失，\(x_1a\ne0\)，且
\(\nu(x_1a)=1+\nu(a)\)。这证明 \(x_1\) 是非零因子，也给出
\[
(x_1)\cap\mathfrak m^n=x_1\mathfrak m^{n-1}\quad(n\ge1).
\]
一边的包含来自 \(x_1\in\mathfrak m\)；反向，若 \(x_1a\in\mathfrak m^n\) 且 \(a\ne0\)，阶数公式给 \(\nu(a)\ge n-1\)，故 \(a\in\mathfrak m^{n-1}\)。这个交式还保证商掉 \(x_1\) 不会产生额外的初始关系。对 \(n\ge1\)，商环伴随分次的第 \(n\) 层为
\[
\dfrac{\mathfrak m^n+(x_1)}{\mathfrak m^{n+1}+(x_1)}
\cong
\dfrac{\mathfrak m^n}{\mathfrak m^{n+1}+x_1\mathfrak m^{n-1}}.
\]
因为 \(\mathfrak m^{n+1}\subseteq\mathfrak m^n\)，分母为 \(\mathfrak m^{n+1}+((x_1)\cap\mathfrak m^n)\)，恰由上述交式得到。右端是 \(k[X_1,\ldots,X_d]/(X_1)\) 的第 \(n\) 次部分；第零层另为 \(R/\mathfrak m=k\)。各层同构都由自然商映射诱导，因而与乘法相容，故
\[
\operatorname{gr}_{\mathfrak m/(x_1)}(R/(x_1))
\cong k[X_1,\ldots,X_d]/(X_1)
\cong k[X_2,\ldots,X_d].
\]
当 \(d\ge2\) 时，在 \(R/(x_1)\) 中，\(x_2\) 的初始式成为剩余多项式环中的变量 \(X_2\)，所以同一论证证明它为非零因子，并把商环的伴随分次环变成 \(k[X_3,\ldots,X_d]\)。逐次重复，每个 \(x_i\) 都在前面的商环上为非零因子，末端商为 \(R/\mathfrak m=k\ne0\)，故整组正则。这个论证适用于任意一组极小生成元。

若反过来 \(\mathfrak m\) 由长度 \(r\) 的正则序列生成，那么
\(r\le\operatorname{depth}R\le\dim R\le\operatorname{edim}R\le r\)。其中深度不超过维数见\cref{b3:cor:depth-dimension-bound}，高度不等式见\eqref{b3:eq:dimension-embedding-bound}。所以这些数均相等，环正则，且深度等于 \(d\)。

任意参数系的断言只需刚得到的 \(\operatorname{depth}R=\dim R=d\)。事实上，以下论证适用于所有满足这个等式的 Noether 局部环，可对 \(d\) 归纳。取参数系 \(a_1,\ldots,a_d\)；\(d=0\) 时为空正则序列。设 \(d\ge1\)。对 \(\mathfrak p\in\operatorname{Ass}_RR\)，\cref{b3:cor:associated-depth-bound}给 \(d\le\dim R/\mathfrak p\le\dim R=d\)，故这些关联素理想的商环维数都恰为 \(d\)。若 \(a_1\in\mathfrak p\)，由 \(\sqrt{(a_1,\ldots,a_d)}=\mathfrak m\) 得
\[
\sqrt{(a_2,\ldots,a_d)(R/\mathfrak p)}
=\mathfrak m/\mathfrak p.
\]
故这 \(d-1\) 个元素在局部环 \(R/\mathfrak p\) 中生成极大理想准素的理想，高度定理使该环维数至多 \(d-1\)，矛盾。因此 \(a_1\) 避开全部关联素理想，由\cref{b3:thm:associated-zero-divisors}是非零因子。\cref{b3:cor:regular-quotient-depth,b3:prop:regular-quotient-dimension}使 \(R/(a_1)\) 的深度和维数均为 \(d-1\)，而根等式在这个商环中仍给出 \(\sqrt{(\bar a_2,\ldots,\bar a_d)}=\mathfrak m/(a_1)\)。余下元素于是构成商环的参数系，归纳得到整组正则。归纳保留的是深度等于维数，不要求中间商环仍然正则。
\end{proof}

\subsection{参数系与正则序列}
\label{b3:subsec:1801:03}

上述结果中的两类参数系在例子中很容易区分。设 \(k\) 为域。在 \(R=k[x,y]_{(x,y)}\) 中，\(x,y\) 是极大理想的极小生成组；\(x^2,y^3\) 也是参数系，且为正则序列，却不能生成极大理想。

\begin{corollary}\label{b3:cor:regular-parameter-quotient}
设 \((R,\mathfrak m)\) 为维数 \(d\) 的交换 Noether 正则局部环，\(k=R/\mathfrak m\)，\(x_1,\ldots,x_d\) 为 \(\mathfrak m\) 的极小生成组。则对整数 \(0\le r\le d\)，
\(R/(x_1,\ldots,x_r)\) 为维数 \(d-r\) 的正则局部环，其极大理想进滤过的伴随分次环与 \(k[X_{r+1},\ldots,X_d]\) 分次同构，其中每个变量次数为一。
\end{corollary}
\begin{proof}
前一定理证明中的逐层交式给出了所述分次同构；Hilbert–Samuel 定理使该商环维数为 \(d-r\)，其极大理想模平方的维数也为 \(d-r\)，故正则。
\end{proof}

若只知道所商掉的元素 \(f\) 是非零因子，商环未必正则。例如在 \(k[x,y]_{(x,y)}\) 中，\(f=x^2\) 是非零因子，但商环的维数为一，极大理想仍需两个生成元。这里 \(f\in\mathfrak m^2\)，已经失去极小生成元的非零一次类。

\begin{proposition}\label{b3:prop:regular-local-basic-examples}
设 \(k\) 为任意域，\(p\) 为素数，记下列环的极大理想分别为 \(\mathfrak m_1,\mathfrak m_2,\mathfrak m_3\)：
\[
R_1=k[x,y,z]_{(x,y,z)}/(z-xy),\qquad
R_2=k[x,y,z]_{(x,y,z)}/(z^2-xy),\qquad
R_3=\mathbb Z[x]_{(p,x)}.
\]
环 \(R_1\) 的维数与嵌入维数均为二，因而正则，且
\[
\operatorname{gr}_{\mathfrak m_1}R_1\cong k[X,Y].
\]
环 \(R_2\) 的维数为二、嵌入维数为三，因而不正则，且
\[
\operatorname{gr}_{\mathfrak m_2}R_2\cong k[X,Y,Z]/(Z^2-XY).
\]
环 \(R_3\) 的维数与嵌入维数均为二，因而正则；其剩余域为 \(\mathbb F_p\)，并有
\[
\operatorname{gr}_{\mathfrak m_3}R_3\cong\mathbb F_p[P,X],
\qquad
\beta_i^{R_3}(\mathbb F_p)=
\begin{cases}
1,&i=0,2,\\
2,&i=1,\\
0,&i>2.
\end{cases}
\]
上述分次同构中的变量均为一次类，\(P\) 对应 \(p+\mathfrak m_3^2\)。
\end{proposition}
\begin{proof}
关系 \(z=xy\) 使 \(k[x,y,z]/(z-xy)\cong k[x,y]\)。极大理想 \((x,y,z)\) 对应 \((x,y)\)，局部化后得到 \(R_1\cong k[x,y]_{(x,y)}\)。因此其维数与嵌入维数都为二，由\cref{b3:thm:regular-local-characterizations}得到所述伴随分次环。这个一次关系消去了 \(z\) 的生成方向。

令 \(A=k[x,y,z]_{(x,y,z)}\)、\(\mathfrak n=(x,y,z)A\)。这是三维局部整环，\(z^2-xy\ne0\)，所以该元素为非零因子；\cref{b3:prop:regular-quotient-dimension}给 \(\dim R_2=2\)。因为 \(z^2-xy\in\mathfrak n^2\)，它在 \(\mathfrak n/\mathfrak n^2\) 中没有给出一次关系，故
\[
\mathfrak m_2/\mathfrak m_2^2
\cong\mathfrak n/(\mathfrak n^2+(z^2-xy))
=\mathfrak n/\mathfrak n^2
\]
仍为三维。这个方程本身齐次，由\cref{b3:prop:hypersurface-tangent-cone}，它的初始关系恰为 \(Z^2-XY\)，得到所述分次商环。这些论证不要求 \(k\) 的特征避开任何素数。

\(\mathbb Z[x]\) 为 Noether 整环，故 \(R_3\) 也是 Noether 局部整环。在其中有严格素理想链
\[
(0)\subsetneq pR_3\subsetneq(p,x)R_3=\mathfrak m_3.
\]
其中 \(R_3/pR_3\cong\mathbb F_p[x]_{(x)}\)，所以前两项为素理想，且 \(x\notin pR_3\)。链给出维数下界二；\(\mathfrak m_3\) 由 \(p,x\) 生成，高度定理给出上界二。因此
\[
2=\dim R_3\le\operatorname{edim}R_3\le2,
\]
环正则，\(p,x\) 的一次类为余切空间的一组基，规范分次同构便给出 \(\mathbb F_p[P,X]\)。乘 \(p\) 在 \(R_3\) 上为单射，商掉 \(p\) 后乘 \(x\) 在 \(\mathbb F_p[x]_{(x)}\) 上也为单射，末端商为 \(\mathbb F_p\ne0\)。故 \(p,x\) 为正则序列，\cref{b3:thm:koszul-regular-resolution}给出秩分别为 \(1,2,1\) 的剩余域自由消解。各微分的系数都属于 \(\mathfrak m_3\)，所以它极小，与 \(\mathbb F_p\) 张量后微分为零，得到所述 Betti 数。
\end{proof}

环 \(\mathbb Z[x]_{(p,x)}\) 本身的特征为零，剩余域的特征却为 \(p\)；它不能含有剩余域作为子域。正则性要求伴随分次环是剩余域上的多项式环，并不要求原环本身是域上的多项式局部环。

% !TeX root = ../../Book3.tex
% 纲要标题：### 18.2 同调刻画
% 纲要位置：计划纲要.md，第 2058 行。

\section{同调刻画}
\label{b3:sec:1802}

正则参数的 Koszul 复形使剩余域具有有限自由消解。反方向更严格：剩余域只要具有有限投射维数，就会迫使环正则，并同时控制所有有限模的消解长度。

% 纲要内容：
%
% * 剩余域的有限投射维数
% * Auslander–Buchsbaum–Serre 正则性判据
% * 正则局部环的整体维数
%

\subsection{剩余域的有限投射维数}
\label{b3:subsec:1802:01}

如果 \(R\) 是 \(d\) 维正则局部环，极大理想的一组极小生成元是正则序列。其 Koszul 复形是剩余域的极小自由消解，故
\begin{equation}\label{b3:eq:regular-residue-betti}
\operatorname{pd}_Rk=d,\qquad \beta_i^R(k)=\binom di.
\end{equation}
真正困难的是反向：自由消解能够终止，为什么会迫使极大理想恰有 \(d\) 个生成元？下面先准备一个选择元素的步骤。

\begin{lemma}\label{b3:lem:linear-regular-element-choice}
设 \((R,\mathfrak m)\) 为交换 Noether 局部环，\(\mathfrak m\ne0\) 且 \(\operatorname{depth}R>0\)。则存在
\(x\in\mathfrak m\setminus\mathfrak m^2\) 在 \(R\) 上为非零因子。
\end{lemma}
\begin{proof}
正深度保证 \(\mathfrak m\notin\operatorname{Ass}R\)。从有限个关联素理想中只保留包含极大的那些，记为 \(\mathfrak p_1,\ldots,\mathfrak p_s\)，它们两两不包含。Nakayama 引理使 \(\mathfrak m\ne\mathfrak m^2\)，取 \(y\in\mathfrak m\setminus\mathfrak m^2\)。

对每个满足 \(y\in\mathfrak p_i\) 的指标，选 \(a_i\in\mathfrak m\setminus\mathfrak p_i\)，并对 \(j\ne i\) 选 \(b_{ij}\in\mathfrak p_j\setminus\mathfrak p_i\)。令
\(w_i=a_i^2\prod_{j\ne i}b_{ij}\)。则 \(w_i\in\mathfrak m^2\)，它不在 \(\mathfrak p_i\) 中，却在其余所有 \(\mathfrak p_j\) 中。于是
\(x=y+\sum_{y\in\mathfrak p_i}w_i\)
模 \(\mathfrak m^2\) 与 \(y\) 相同，并避开每个 \(\mathfrak p_i\)：若 \(y\in\mathfrak p_i\)，唯一不属于该素理想的加项为 \(w_i\)；否则全部修正项均属于它。关联素理想检验使 \(x\) 为非零因子。这个构造不要求剩余域无限。
\end{proof}

\subsection{Auslander–Buchsbaum–Serre 正则性判据}
\label{b3:subsec:1802:02}

反向证明要把“剩余域的消解有限”变成“极大理想的生成元数等于维数”。先用 Auslander–Buchsbaum 公式取得正深度，选一个既不在 \(\mathfrak m^2\) 中又不是零因子的 \(x\)；再把 \(\mathfrak m/x\mathfrak m\) 的有限投射维数传给商环极大理想 \(\mathfrak m/xR\)。这两个模并不相同，下面的直和分裂正是传递有限性的关键。如此每次降低一个嵌入维数，归纳得到正则性；从有限模推广到任意模则是另一步，不能仅靠模的有向并完成。

\begin{theorem}[Auslander–Buchsbaum–Serre]\label{b3:thm:regular-local-global-dimension}
设 \((R,\mathfrak m,k)\) 为交换 Noether 局部环。则 \(R\) 正则，当且仅当 \(\operatorname{pd}_Rk<\infty\)。此时每个 \(R\) 模的投射维数至多 \(\dim R\)，并且
\[
\operatorname{gldim}R=\operatorname{pd}_Rk=\dim R,
\qquad
\operatorname{gldim}R\coloneqq\sup_{M}\operatorname{pd}_RM,
\]
其中上确界取遍全部 \(R\) 模，包括非有限生成模。
\end{theorem}

这称为\term[auslander-buchsbaum-serre-panju]{Auslander–Buchsbaum–Serre 正则性判据}[Auslander--Buchsbaum--Serre criterion]；\(\operatorname{gldim}R\) 称为环的\term[zhengti-weishu]{整体维数}[global dimension]。下面先证明等价性，随后补上从有限模到所有模的步骤。这一证明依照\cite[Theorem 19.2, pp. 156--157]{Matsumura1989CommutativeRingTheory}中对极大理想的处理。

\begin{proof}[\cref{b3:thm:regular-local-global-dimension}的证明]
正则性推出 \(\operatorname{pd}_Rk=d\) 已由\eqref{b3:eq:regular-residue-betti}得到。反过来假设 \(\operatorname{pd}_Rk<\infty\)，对嵌入维数 \(e\) 归纳。\(e=0\) 时 Nakayama 引理给 \(\mathfrak m=0\)，环为域。

设 \(e>0\)。若 \(k\) 投射，则它有限自由；模 \(\mathfrak m\) 后比较秩可知秩为一，故 \(\operatorname{Ann}_Rk=0\)，从而 \(\mathfrak m=0\)，矛盾。因此 \(\operatorname{pd}_Rk>0\)，Auslander–Buchsbaum 公式及 \(\operatorname{depth}_Rk=0\) 给 \(\operatorname{depth}R>0\)。由\cref{b3:lem:linear-regular-element-choice}取正则元 \(x\in\mathfrak m\setminus\mathfrak m^2\)，令 \(B=R/(x)\)、\(\mathfrak n=\mathfrak m/(x)\)。

短正合列 \(0\to\mathfrak m\to R\to k\to0\) 说明 \(\mathfrak m\) 投射维数有限。\(x\) 在 \(\mathfrak m\subseteq R\) 上也为非零因子，故\cref{b3:lem:regular-element-resolution-basechange}使 \(\mathfrak m/x\mathfrak m\) 成为有限投射维数的 \(B\) 模。

必须注意 \(\mathfrak m/x\mathfrak m\) 并不等于 \(\mathfrak n\)。把 \(x\) 扩充为极小生成组 \(x,x_2,\ldots,x_e\)，令 \(\mathfrak b=(x_2,\ldots,x_e)\)。若 \(xa\in\mathfrak b\) 且 \(a\) 为单位，则 \(x\in\mathfrak b\)，违反极小性。因此
\(\mathfrak b\cap xR\subseteq x\mathfrak m\)，从而自然映射
\[
\mathfrak n\cong\mathfrak b/(\mathfrak b\cap xR)
\longrightarrow\mathfrak m/x\mathfrak m
\longrightarrow\mathfrak m/xR=\mathfrak n
\]
的复合为恒等。\(\mathfrak n\) 是直和因子，故由\cref{b3:cor:residue-pd-bound}的直和因子结论，它在 \(B\) 上投射维数有限。给 \(\mathfrak n\) 的有限自由消解末端接上
\(0\to\mathfrak n\to B\to k\to0\)，便知 \(\operatorname{pd}_Bk<\infty\)。

\(B\) 的嵌入维数为 \(e-1\)，由归纳它正则。又因 \(x\) 为非零因子，\cref{b3:prop:regular-quotient-dimension}给
\(\dim R=\dim B+1=e\)。所以 \(R\) 正则。由\cref{b3:cor:residue-pd-bound}，所有有限 \(R\) 模的投射维数至多 \(d\)。下一小节的\cref{b3:lem:cyclic-global-dimension-test}把此界推广到所有模，得到 \(\operatorname{gldim}R\le d\)；取 \(M=k\) 得反向不等式。
\end{proof}

\subsection{正则局部环的整体维数}
\label{b3:subsec:1802:03}

有限模上的统一界并不靠“每个模都是有限模的并”直接推广；自由性和投射性不由任意有向并保持。需要使用内射对象检验 Ext。

\begin{lemma}\label{b3:lem:cyclic-global-dimension-test}
设 \(R\) 为交换环，\(n\ge0\)。若对每个理想 \(I\subseteq R\)，有 \(\operatorname{pd}_R(R/I)\le n\)，则对每个 \(R\) 模 \(M\)，均有 \(\operatorname{pd}_RM\le n\)。
\end{lemma}
\begin{proof}
这里使用第二卷第7.2节已经证明的 Baer 判据和足够多内射模定理：每个模能嵌入内射模，且模 \(E\) 内射恰好等价于每个 \(I\to E\) 能延拓到 \(R\)。后一条件由 \(0\to I\to R\to R/I\to0\) 的 Hom 长正合列，等价于 \(\operatorname{Ext}_R^1(R/I,E)=0\) 对全部理想成立。

固定任意模 \(N\)，连续把它及所得余核嵌入内射模，得到
\(0\to N\to E^0\to\cdots\to E^{n-1}\to C\to0\)；\(n=0\) 时取 \(C=N\)。内射模使正次 \(\operatorname{Ext}_R^i(M,E^j)\) 为零，因为 \(\operatorname{Hom}_R(-,E^j)\) 正合。沿第二变量的长正合列平移，得到
\[
\operatorname{Ext}_R^1(R/I,C)
\cong\operatorname{Ext}_R^{n+1}(R/I,N)=0.
\]
Baer 判据使 \(C\) 内射。因此对任意 \(M\)，同样平移给出
\(\operatorname{Ext}_R^{n+1}(M,N)\cong\operatorname{Ext}_R^1(M,C)=0\)。这对任意 \(N\) 成立。

取 \(M\) 的自由消解，截到第 \(n\) 个合冲模 \(K\)。第一变量的维数平移给 \(\operatorname{Ext}_R^1(K,N)=0\) 对所有 \(N\) 成立；\secref{b3:sec:1702}中已说明，这使到 \(K\) 的自由满射分裂，故 \(K\) 投射。于是截断的序列是长度 \(n\) 的投射消解。
\end{proof}

因此前一定理对有限模得到的维数界也适用于任意模。整体维数的定义取遍所有模；在 Noether 环上，有限模给出同一上确界，依据正是上述 Baer 判据与维数平移的论证。

% !TeX root = ../../Book3.tex
% 纲要标题：### 18.3 局部化与完备化
% 纲要位置：计划纲要.md，第 2064 行。

\section{局部化与完备化}
\label{b3:sec:1803}

正则性应能在缩小到素点附近或转入形式近似后继续使用。局部化通过有限消解保持这一性质，完备化则保留维数与余切空间；两种操作的证明依据不同。

% 纲要内容：
%
% * 正则性在局部化下的行为
% * 正则性与完备化
% * 正则局部环是整环的证明方法
%

\subsection{正则性的局部化}
\label{b3:subsec:1803:01}

\begin{theorem}\label{b3:thm:regular-local-localization}
设 \(R\) 为交换 Noether 正则局部环，\(\mathfrak p\in\operatorname{Spec}R\)。则 \(R_{\mathfrak p}\) 仍是正则局部环。
\end{theorem}
\begin{proof}
有限 \(R\) 模 \(R/\mathfrak p\) 由\cref{b3:thm:regular-local-global-dimension}具有有限自由消解。局部化正合，故该消解在 \(\mathfrak p\) 处给出
\((R/\mathfrak p)_{\mathfrak p}=\kappa(\mathfrak p)\) 的有限自由 \(R_{\mathfrak p}\) 消解。后者是 Noether 局部环的剩余域，再用同一定理的反向便得正则性。
\end{proof}

\begin{definition}\label{b3:def:regular-ring}
设 \(R\) 为交换 Noether 环。若每个素理想 \(\mathfrak p\) 的局部环 \(R_{\mathfrak p}\) 都正则，则称 \(R\) 为\term[zhengzehuan]{正则环}[regular ring]。
\end{definition}

只检验极大理想已足够：每个素理想包含在某个极大理想中，再使用前一定理。正则环未必局部，也未必整环，例如两个域的直积。关于正则局部环的断言若去掉“局部”，需要重新检查结论。

\subsection{正则性与完备化}
\label{b3:subsec:1803:02}

\begin{theorem}\label{b3:thm:regular-local-completion}
设 \((R,\mathfrak m,k)\) 为交换 Noether 局部环，\(\widehat R\) 为其 \(\mathfrak m\) 进完备化。则 \(R\) 正则当且仅当 \(\widehat R\) 正则。
\end{theorem}
\begin{proof}
由\cref{b3:thm:completion-noetherian}，\(\widehat R\) 是 Noether 局部环，极大理想为 \(\mathfrak m\widehat R\)，剩余域仍为 \(k\)。\cref{b3:prop:completion-quotients}在二次幂处给出
\[
\mathfrak m/\mathfrak m^2
\cong\mathfrak m\widehat R/(\mathfrak m\widehat R)^2,
\]
所以两环嵌入维数相同。\cref{b3:thm:completion-dimension}使它们 Krull 维数相同。定义中的两个数分别不变，故正则性等价。
\end{proof}

例如 \(k[x_1,\ldots,x_d]_{(x_1,\ldots,x_d)}\) 的完备化为 \(k[\![x_1,\ldots,x_d]\!]\)，所以后者正则。这里保留全部形式幂级数，不能把它与只允许有限项的多项式局部环当成同一个环。

\subsection{正则局部环的整性}
\label{b3:subsec:1803:03}

\begin{proposition}\label{b3:prop:regular-local-domain}
交换 Noether 正则局部环是整环。
\end{proposition}
\begin{proof}
设极大理想为 \(\mathfrak m\)。若 \(a,b\ne0\)，Krull 交定理使它们具有有限阶数 \(r,s\)，故在伴随分次环中给出非零初始式。\cref{b3:thm:regular-local-characterizations}说明该分次环是域上的多项式环，因而两初始式的乘积非零。这一乘积正是 \(ab\) 在 \(\mathfrak m^{r+s}/\mathfrak m^{r+s+1}\) 中的像，所以 \(ab\ne0\)。
\end{proof}

这个证明也说明分离性和分次结构各自的作用：前者保证非零元素有初始式，后者保证相乘时最低项不会消失。一般环的伴随分次环不一定是整环，即使原环是整环；尖点的切锥便给出反例。

\ProseParagraphBreak
一维正则局部环的极大理想由一个元素生成，再结合刚得到的整性，它由\cref{b3:thm:dvr-characterizations}为 DVR。下一节将把唯一分解推广到任意有限维数；这需要比“一维局部化是 DVR”更多的信息。

\ProseParagraphBreak
到此为止，正则性已有三种用法：极小生成元给出正则序列，剩余域的有限消解检测正则性，局部化与完备化把它带到新的局部环。阶乘性还要解决另一个问题：高度一素理想在各点附近可由一元生成，是否存在同一个全局生成元。下一节用行列式线回答这一问题；局部化保持正则性本身尚不能回答它。

% !TeX root = ../../Book3.tex
% 纲要标题：### 18.4 正则局部环的阶乘性
% 纲要位置：计划纲要.md，第 2070 行。

\section{正则局部环的阶乘性}
\label{b3:sec:1804}

唯一分解要求高度一素理想由单个元素生成。正则局部环的有限自由消解和行列式线能够证明这一主性，而更高高度的理想一般仍须多个生成元。

% 纲要内容：
%
% * 证明正则局部环是唯一分解整环，先把目标化为高度一素理想的主性；完整保留 Noether 局部假设
% * 准备所需行列式线引理：利用有限自由消解控制主开集上的可逆模；在使用处明确有限展示模延拓与局部化正合性的论证
% * 选择 \(x\in\mathfrak m\setminus\mathfrak m^2\)，结合 \(R/(x)\) 的正则性、维数归纳及主开集上的秩一模分析，完成高度一素理想主性的证明
% * 区分“每个素点附近局部为主”与“给定主开集上的可逆模为自由”两个步骤；不能只凭局部阶乘性推断全局唯一分解
% * 将第一卷唯一分解、第二卷行列式线与本卷局部同调结构连成完整证明链；第20章据此讨论除子与局部阶乘性
%

\subsection{唯一分解与高度一素理想的主性}
\label{b3:subsec:1804:01}

第一卷的唯一分解理论把元素的分解归结为不可约元素是否为素元。在 Noether 整环中，还可以把目标改写为高度一素理想的主性。这一形式便于在局部化中比较。

\begin{proposition}\label{b3:prop:height-one-ufd-criterion}
交换 Noether 整环 \(A\) 是唯一分解整环，当且仅当每个高度一素理想都是主理想。
\end{proposition}
\begin{proof}
若 \(A\) 是唯一分解整环，取高度一素理想 \(\mathfrak p\) 中的非零元并将它分解为素元之积。至少一个素元 \(a\) 属于 \(\mathfrak p\)，于是
\((0)\subsetneq(a)\subseteq\mathfrak p\)。高度一迫使 \(\mathfrak p=(a)\)。

反过来，Noether 性保证每个非零非单位可分解为有限个不可约元素：否则反复选取仍不能分解的真因子，得到严格递增的主理想链，违反升链条件。取不可约元 \(a\)，并取 \((a)\) 上的极小素理想 \(\mathfrak p\)。主理想定理及整性给 \(\operatorname{ht}\mathfrak p=1\)，故 \(\mathfrak p=(b)\)。写 \(a=bc\)，其中 \(b\) 非单位，由不可约性知 \(c\) 为单位。所以 \((a)=\mathfrak p\)，\(a\) 为素元。所有不可约元素都是素元，分解存在且唯一到次序与单位。
\end{proof}

判据及下面的下降步骤可对照\cite[Theorems 20.1--20.2, pp. 161--162]{Matsumura1989CommutativeRingTheory}。高度一的局部计算还不能直接给出全局生成元；以下行列式论证专门处理这个问题。

证明分三步。对高度一素理想 \(\mathfrak p\) 选正则参数 \(x\)：若 \(x\in\mathfrak p\)，则 \(\mathfrak p=(x)\)；否则在主开集 \(D(x)\) 上，用较低维局部环的阶乘性说明 \(\mathfrak pR_x\) 逐点自由秩一。有限展示把这种逐点自由性提升为可逆模，行列式线再利用原环上的有限自由消解给它一个全局基；最后利用 \(x\) 是素元，把 \(R_x\) 中的主生成元下降回 \(R\)。\cref{b3:thm:regular-local-ufd}将完成这一证明；以下小节先建立其中新增的两个工具。

\subsection{行列式线与主开集上的可逆模}
\label{b3:subsec:1804:02}

第二卷第9.2—9.3节已经建立外幂、直和的外幂分解和行列式线。对交换整环 \(A\) 上常秩 \(r\) 的有限投射模 \(P\)，记
\(\det P=\bigwedge_A^rP\)，并约定 \(\det0=A\)。它是秩一投射模；局部取基即可验证这一点。秩一有限投射模也称可逆模，其对偶是张量逆。

\begin{lemma}\label{b3:lem:determinant-finite-free-resolution}
设 \(A\) 为交换整环，\(L\) 为秩一有限投射 \(A\) 模。若 \(L\) 具有有限自由消解，则 \(L\cong A\)。
\end{lemma}
\begin{proof}
设消解为 \(0\to F_s\to\cdots\to F_0\to L\to0\)，令 \(K_0=L\)，并把它分成
\(0\to K_{i+1}\to F_i\to K_i\to0\)。由于 \(K_0\) 投射，第一列分裂，\(K_1\) 是有限自由模的直和项，故有限投射。依次进行，所有 \(K_i\) 有限投射，最后 \(K_s\cong F_s\)。其秩为常数：局部自由秩都等于张量到分式域后的维数。

对分裂短正合列 \(0\to P\to F\to Q\to0\)，取一个截面得 \(F\cong P\oplus Q\)。第二卷的直和外幂公式在最高次给出
\(\det F\cong\det P\otimes_A\det Q\)。这个同构不依赖截面：局部以 \(P\) 的基开头，再接 \(Q\) 的基的提升；改变提升只是给后面的列加上前面列的线性组合，最高楔积不变。局部同构因而相容，给出全局同构。

把各短正合列的行列式等式依次代入，得到
\[
L=\det L\cong
\det F_0\otimes(\det F_1)^*\otimes\det F_2
\otimes\cdots\otimes(\det F_s)^{(-1)^s},
\]
其中负指数表示取对偶。每个 \(F_i\) 自由，其最高外幂也自由秩一，右边因而同构于 \(A\)。
\end{proof}

\begin{lemma}\label{b3:lem:principal-open-invertible-free}
设 \(R\) 为交换 Noether 正则局部环，\(0\ne f\in R\)。每个可逆 \(R_f\) 模都自由秩一；换言之，\(\operatorname{Pic}(R_f)=0\)。
\end{lemma}
\begin{proof}
给定可逆模 \(L\)，它有限投射，故有限展示。选展示矩阵
\(D:R_f^a\to R_f^b\)，使 \(L\cong\operatorname{coker}D\)。只有有限个矩阵元素，可取一个共同分母 \(f^N\)，使 \(D'=f^ND\) 的元素全在 \(R\) 中。令
\(M=\operatorname{coker}(D':R^a\to R^b)\)。局部化后 \(f^N\) 为单位，两个矩阵有相同像，故 \(M_f\cong L\)。这说明延拓的是一个有限展示模，不是先假定 \(L\) 能延拓为自由或投射模。

由\cref{b3:thm:regular-local-global-dimension}，有限生成模 \(M\) 有有限投射维数；其极小消解各项有限自由，故是有限自由消解。局部化保持正合，把它变为 \(L\) 的有限自由 \(R_f\) 消解。\(R_f\) 是整环，由\cref{b3:lem:determinant-finite-free-resolution}，\(L\cong R_f\)。
\end{proof}

\subsection{正则参数、维数归纳与秩一模}
\label{b3:subsec:1804:03}

现在取 \(d\) 维正则局部环 \((R,\mathfrak m)\)、\(d>0\)，选
\(x\in\mathfrak m\setminus\mathfrak m^2\)。由\cref{b3:cor:regular-parameter-quotient}，\(R/(x)\) 正则，因而由\cref{b3:prop:regular-local-domain}为整环。所以 \(x\) 不仅是非零因子，还是素元。

\ProseParagraphBreak
假设较低维的正则局部环都已知为唯一分解整环。对任意非极大素理想 \(\mathfrak q\)，局部环 \(R_{\mathfrak q}\) 正则，而且维数小于 \(d\)：其素链在 \(R\) 中末端还可接上 \(\mathfrak m\)。因此维数归纳适用于所有这样的点。

\begin{lemma}\label{b3:lem:prime-generator-descends}
设 \(R\) 为交换 Noether 整环，\(x\) 为素元，\(\mathfrak p\) 为不含 \(x\) 的素理想。若 \(\mathfrak pR_x\) 主生成，则 \(\mathfrak p\) 主生成。
\end{lemma}
\begin{proof}
把局部生成元乘一个 \(x\) 的幂，可写成 \(f\in\mathfrak p\)，且 \(\mathfrak pR_x=fR_x\)。若 \(x\mid f\)，约去一次 \(x\) 后的元仍属于 \(\mathfrak p\)，因为它素且不含 \(x\)。这个约去过程不能无限继续，否则得到严格递增的主理想链。故可选 \(f\notin xR\)。

对 \(z\in\mathfrak p\)，局部等式给出 \(x^nz=fh\)，其中 \(h\in R\)、\(n\ge0\)。若 \(n>0\)，\(x\) 素且不整除 \(f\)，故 \(x\mid h\)。约去一个 \(x\) 并重复，得到 \(h=x^nh'\)，于是 \(z=fh'\)。这证明 \(\mathfrak p\subseteq fR\)，反向包含由 \(f\in\mathfrak p\) 显然。
\end{proof}

\subsection{局部主性与主开集上的自由性}
\label{b3:subsec:1804:04}

设 \(\mathfrak p\) 为高度一素理想且 \(x\notin\mathfrak p\)。在 \(R_x\) 的任意素点 \(\mathfrak qR_x\) 处，\(\mathfrak q\ne\mathfrak m\)。若 \(\mathfrak p\nsubseteq\mathfrak q\)，局部化后的理想是整个 \(R_{\mathfrak q}\)；若 \(\mathfrak p\subseteq\mathfrak q\)，它仍是高度一素理想，较低维的唯一分解性使它主生成。因此 \(\mathfrak pR_x\) 在每个素点都自由秩一。

\ProseParagraphBreak
这里有两次不同的推论。首先，\(R_x\) Noether 使 \(\mathfrak pR_x\) 有限展示，\cref{b3:thm:finite-presentation-flat-projective}才允许从各点自由推出它是可逆模。其次，\cref{b3:lem:principal-open-invertible-free}利用原环 \(R\) 的有限自由消解，才把这个可逆模变成 \(R_x\) 上的自由模。前一步本身没有产生全局基。

\ProseParagraphBreak
第二步不能从论证中省去。\chapref{b3:ch:10}中 \(\mathbb Z[\sqrt{-5}]\) 的理想 \((2,1+\sqrt{-5})\) 在各素点局部主生成，却不是主理想。局部主性的结论在该例中完全成立，缺少的正是保证可逆模自由的进一步条件。

\subsection{正则局部环阶乘性的证明及其与除子的联系}
\label{b3:subsec:1804:05}

\begin{theorem}\label{b3:thm:regular-local-ufd}
每个交换 Noether 正则局部环都是唯一分解整环。
\end{theorem}
\begin{proof}
由\cref{b3:prop:regular-local-domain}，它首先是整环。对维数 \(d\) 归纳。\(d=0\) 时为域。设 \(d>0\)，取 \(x\in\mathfrak m\setminus\mathfrak m^2\)。\cref{b3:cor:regular-parameter-quotient}及\cref{b3:prop:regular-local-domain}使 \(x\) 为素元。

取高度一素理想 \(\mathfrak p\)。若 \(x\in\mathfrak p\)，则非零素理想 \((x)\subseteq\mathfrak p\) 由高度一得到相等。若 \(x\notin\mathfrak p\)，对每个 \(R_x\) 素点，\cref{b3:thm:regular-local-localization}及维数归纳使其局部环为唯一分解整环。因此上一小节的两步论证适用：\(\mathfrak pR_x\) 是可逆模，继而由\cref{b3:lem:principal-open-invertible-free}为自由秩一，故主生成。\cref{b3:lem:prime-generator-descends}再使 \(\mathfrak p\) 在 \(R\) 中主生成。

所有高度一素理想都已主生成，\cref{b3:prop:height-one-ufd-criterion}给出唯一分解性。
\end{proof}

这一同调证明的经典形式见\cite[Theorem 20.3, pp. 163--164]{Matsumura1989CommutativeRingTheory}。它把自由消解用于控制秩一模，进而回到元素分解。到\chapref{b3:ch:20}讨论除子时，高度一素理想将成为除子的基本组成部分；本节已经证明，在正则局部环中这些组成部分都由一个元素给出。

% !TeX root = ../../Book3.tex
% 纲要标题：### 18.5 正则点、切空间与光滑性
% 纲要位置：工作记录/计划纲要.md，第 2580 行。
% 纲要细目（写作范围）：
% * 正则点以局部环余切空间检验，与相对微分纤维通过剩余域余切序列比较；在有理超曲面点写出维数与嵌入维数
% * 对照环论正则性、基域上的光滑性与几何正则性；有限型代数的光滑与几何正则比较引用 Stacks Project Lemma 33.12.6
% * 以纯不可分域扩张说明正则不必几何正则；第23章证明完美基域上正则与光滑的逐点比较
% * 正文命题计算不可分闭点的余切序列，说明局部余切空间与微分纤维同维仍不保证自然映射单射
% ---

\section{正则点、切空间与光滑性}
\label{b3:sec:1805}

几何中的点由相应局部环描述，正则性比较该环的维数与一阶参数数目。相对光滑性还依赖基域；在不完美域上，即使环自身正则，改变基域后也可能出现幂零元。

% 纲要内容（仅作写作计划，尚非教材正文）：
%
% * 正则点与切空间维数
% * 对照环论正则性、基域上的光滑性与几何正则性；有限型代数的光滑与几何正则比较引用 Stacks Project Lemma 33.12.6
% * 以纯不可分域扩张说明正则不必几何正则；第23章证明完美基域上正则与光滑的逐点比较
%
% ---
%

\subsection{正则点与切空间维数}
\label{b3:subsec:1805:01}

设 \(A\) 为域 \(k\) 上的有限型代数，\(\mathfrak p\in\operatorname{Spec}A\)。若 \(A_{\mathfrak p}\) 正则，就称 \(\mathfrak p\) 为\term[zhengzedian]{正则点}[regular point]。判定用的是该局部环的剩余域 \(\kappa(\mathfrak p)\)，即比较
\[
\dim A_{\mathfrak p}\quad\text{与}\quad
\dim_{\kappa(\mathfrak p)}
\mathfrak pA_{\mathfrak p}/(\mathfrak pA_{\mathfrak p})^2.
\]
这个商空间是局部环的余切空间，其对偶为该点的 Zariski 切空间。它与相对于 \(k\) 的微分模纤维不是同一个定义。在 \(k\) 有理闭点，剩余域就是 \(k\)，\cref{b3:prop:tangent-cotangent-rational}给出两者的自然同构；一般点则有余切序列
\[
\mathfrak pA_{\mathfrak p}/(\mathfrak pA_{\mathfrak p})^2
\longrightarrow
\kappa(\mathfrak p)\otimes_A\Omega_{A/k}
\longrightarrow
\Omega_{\kappa(\mathfrak p)/k}\longrightarrow0.
\]
末项可能非零，第一箭头也未必单射。因此不能只计算中项的维数，就把它当成局部环的嵌入维数。

\begin{proposition}\label{b3:prop:regular-hypersurface-rational-point}
设 \(k\) 为域，\(n\ge1\) 为整数，\(0\ne f\in k[X_1,\ldots,X_n]\)，且 \(a\in k^n\) 满足 \(f(a)=0\)。则超曲面 \(f=0\) 在 \(a\) 处的局部环正则，当且仅当至少一个 \(\partial f/\partial X_i(a)\ne0\)。
\end{proposition}
\begin{proof}
令 \(S=k[X_1,\ldots,X_n]_{(X_1-a_1,\ldots,X_n-a_n)}\)，其极大理想记为 \(\mathfrak n\)。\(S\) 正则且为整环，\(f\) 是非零因子，因此 \(S/(f)\) 的维数为 \(n-1\)。其极大理想模平方为
\(\mathfrak n/(\mathfrak n^2+(f))\)。\(f\) 在 \(\mathfrak n/\mathfrak n^2\) 中的像由一次 Taylor 项给出，系数恰为各偏导数在 \(a\) 处的值。因此
\[
\dim S/(f)=n-1,\qquad
\operatorname{edim}S/(f)=
\begin{cases}
n-1,&\text{至少一个偏导数在}\;a\;\text{处非零},\\
n,&\text{所有偏导数在}\;a\;\text{处为零}.
\end{cases}
\]
两者相等恰好给出所述条件。
\end{proof}

例如 \(y^2-x^3=0\) 的偏导数为 \(-3x^2,2y\)，在原点总是全零，包括特征二、三的情形；等价地，方程没有一次项。因此原点的局部环维数为一、嵌入维数为二，在任何特征都不正则。\(y-x^2=0\) 对 \(y\) 的偏导数为一，所以每个有理点正则。这一判断检测局部结构，并不依赖图形看起来是否有“尖角”。

\subsection{正则性与基域上的光滑性}
\label{b3:subsec:1805:02}

正则性是环自身的性质。光滑性则是基环到该环的同态的性质，\cref{b3:def:finite-presentation-infinitesimal-properties}用有限展示及平方零提升的存在性定义它。两者所问的问题不同，因此谈论光滑时必须注明基域或基环。

\ProseParagraphBreak
在前一命题的有理点处，如果某个偏导数不为零，可在一个主开邻域上使它可逆。对平方零提升问题，先任意提升各坐标，再用该偏导数的逆修正一个坐标，消掉方程值。这正是\secref{b3:sec:1405}的 Jacobian 线性修正，因而该主开邻域对 \(k\) 光滑。要把这种比较推广到任意点，还须考察扩张基域以后的正则性。

\begin{definition}\label{b3:def:geometrically-regular-algebra}
设 \(k\) 为域，\(A\) 为有限型交换 \(k\) 代数。若对每个域扩张 \(L/k\)，\(A\otimes_kL\) 都是正则环，则称 \(A\) 在 \(k\) 上\term[jihezhengze]{几何正则}[geometrically regular]。
\end{definition}

有限型性保证这些基变换后的环仍 Noether。几何正则性不仅要求原环正则，还要求这一性质在扩张基域后保留。三个条件的区别列于\cref{b3:tab:regular-smooth-geometric}。

\begin{center}
\begin{minipage}{\textwidth}
\makeatletter
\def\@captype{table}
\makeatother
\centering
\caption[正则、光滑与几何正则]{正则、光滑与几何正则。设 \(k\) 为域，\(A\) 为有限型交换 \(k\) 代数；第一行的条件对每个 \(\mathfrak p\in\operatorname{Spec}A\) 检验，\(\mathfrak m_{\mathfrak p}=\mathfrak pA_{\mathfrak p}\)，\(\kappa(\mathfrak p)\) 为其剩余域。}
\label{b3:tab:regular-smooth-geometric}
\begin{tabularx}{\textwidth}{@{}>{\raggedright\arraybackslash}p{0.23\textwidth}>{\raggedright\arraybackslash}X@{}}
\toprule
性质 & 检验内容\\
\midrule
\(A\) 正则 & 每个局部环满足
\(\dim A_{\mathfrak p}=\dim_{\kappa(\mathfrak p)}\mathfrak m_{\mathfrak p}/\mathfrak m_{\mathfrak p}^2\)；条件只涉及 \(A\) 自身。\\
\(A\) 在 \(k\) 上光滑 & \(A\) 有限展示，且从 \(A\) 到交换 \(k\) 代数的平方零商的同态总能提升；条件涉及给定的 \(k\) 代数结构。\\
\(A\) 在 \(k\) 上几何正则 & 对每个域扩张 \(L/k\)，\(A\otimes_kL\) 的所有局部环仍正则；条件检验扩张基域后的结构。\\
\bottomrule
\end{tabularx}
\end{minipage}
\end{center}

\begin{claim}[光滑性与几何正则性]\label{b3:claim:smooth-geometrically-regular}
设 \(k\) 为域，\(A\) 为有限型交换 \(k\) 代数。则 \(A\) 在 \(k\) 上光滑，当且仅当它在 \(k\) 上几何正则。
\end{claim}
这里把一般基域上的比较作为外部结论引用，其完整证明见 Stacks Project, Lemma 33.12.6（\url{https://stacks.math.columbia.edu/tag/038X}），涉及一般光滑性理论和基变换。第23章的\cref{b3:thm:perfect-field-regular-smooth}将证明完美基域上有限型代数的逐点比较：局部环正则当且仅当该点附近对基域光滑。这个后续证明不代替上述任意基域上的一般结论。下面的例子说明，不完美基域上不能省去基变换的检验。

\ProseParagraphBreak
已经得到的有理点判据不能直接照搬到剩余域不可分的点。\eqref{b3:eq:cotangent-residue-sequence}中的剩余域微分项此时可能非零，\(\Omega_{A/k}\) 的纤维不再只记录极大理想的一次部分。下面把这一差别转化为基域扩张后的局部环计算。

\subsection{非完美域上的区分例子}
\label{b3:subsec:1805:03}

设 \(p\) 为素数，\(s\) 为不定元，取 \(k=\mathbb F_p(s)\)、\(K=k[u]/(u^p-s)\)。\cref{b3:exm:inseparable-tangent}已经说明 \(K\) 是域，所以作为局部环，它的极大理想为零，维数和嵌入维数都为零，因而正则。但它对 \(k\) 不光滑：\secref{b3:sec:1405}的平方零提升计算使用 \(C=k[T]/((T^p-s)^2)\)，其平方零理想由 \(T^p-s\) 的像生成。关系的导数为零，使坐标的任何平方零修正都不能消去非零误差 \(T^p-s\)；因此 \(K\to C/(T^p-s)C\) 的恒等映射不能提升。

\ProseParagraphBreak
现在把基域扩张到 \(K\)。有显式同构
\[
K\otimes_kK\cong K[T]/(T^p-s)
\cong K[v]/(v^p),\qquad v=T-u.
\]
最后一个环的唯一素理想为 \((v)\)，所以维数为零；而 \((v)/(v^2)\) 在 \(K\) 上维数为一，所以它不正则。因此 \(K\) 在 \(k\) 上也不几何正则。这里环 \(K\) 自身完全没有零因子，基变换却产生幂零元。这说明“正则局部环”不包含相对于任意给定基域的可分性信息。

\ProseParagraphBreak
相反，多项式代数 \(k[x_1,\ldots,x_d]\) 的平方零提升只需任意提升各坐标像，因而对 \(k\) 光滑；扩张任何基域后仍是多项式代数，故由上述比较几何正则。这并不要求它的每个闭点剩余域都可分于 \(k\)。下面的不可分闭点就位于光滑的仿射直线上，其微分纤维却不能用有理点的公式解释。

\begin{proposition}\label{b3:prop:inseparable-cotangent-fiber}
设 \(p\) 为素数，\(s\) 为 \(\mathbb F_p\) 上的不定元，令
\[
k=\mathbb F_p(s),\qquad A=k[T],\qquad
f=T^p-s,\qquad\mathfrak p=(f),\qquad K=A/\mathfrak p.
\]
则 \(A_{\mathfrak p}\) 是一维正则局部环。在余切序列
\[
\mathfrak pA_{\mathfrak p}/(\mathfrak pA_{\mathfrak p})^2
\longrightarrow \Omega_{A/k}\otimes_AK
\longrightarrow\Omega_{K/k}\longrightarrow0
\]
中，三项分别为以 \(f,\mathrm{d}T,\mathrm{d}\bar T\) 的像为基的一维 \(K\) 向量空间，第一箭头为零，第二箭头为同构。
\end{proposition}
\begin{proof}
\cref{b3:exm:inseparable-tangent}证明了 \(f\) 不可约，所以 \(\mathfrak p\) 是非零极大理想。主理想整环 \(k[T]\) 在 \(\mathfrak p\) 处的局部化是 DVR，其极大理想由 \(f\) 生成，维数和嵌入维数均为一。乘 \(f\) 给出
\[
A_{\mathfrak p}/fA_{\mathfrak p}
\xrightarrow{\sim}
fA_{\mathfrak p}/f^2A_{\mathfrak p};
\]
它单射是因为 \(A_{\mathfrak p}\) 为整环，故左项就是以 \(f\) 的像为基的 \(K\) 向量空间。

多项式微分模为 \(A\,\mathrm{d}T\)，而商代数的微分展示为
\(\Omega_{K/k}=K\,\mathrm{d}\bar T/(f'(\bar T)\,\mathrm{d}\bar T)\)。
特征 \(p\) 下 \(f'=0\)，所以末项也是一维。余切映射把 \(f\) 的像送到 \(\mathrm{d}f=0\)，其后把 \(\mathrm{d}T\) 送到 \(\mathrm{d}\bar T\)，得到所述映射。这里局部余切空间和相对微分纤维虽恰好同维，自然映射却为零；维数相同不足以恢复有理点处的同构。
\end{proof}


\begin{chaptersummary}
正则局部环的三种描述分别体现生成元数、初始结构和消解长度。Hilbert–Samuel 增长排除了伴随分次环中的额外关系，正则参数的 Koszul 复形给出剩余域的有限消解；反向则通过选择一次非零的正则元，并分裂 \(\mathfrak m/x\mathfrak m\to\mathfrak m/xR\)，逐步降低嵌入维数。

\ProseParagraphBreak
正则性在局部化和完备化中保持，但唯一分解还需要单独证明。维数归纳先给出高度一素理想在主开集上的局部主性，有限展示把它变成可逆模；有限自由消解的行列式线再给出全局基，最后沿素元下降到原环。这个中间步骤说明，为何一般的局部主理想并不自动主生成。几何上，正则性比较局部维数与一次参数数目，几何正则性还检验基域扩张后的局部环。纯不可分域扩张的例子显示，原环正则并不足以保证相对光滑性。
\end{chaptersummary}

\BookExercises

\begin{exercise}\label{b3:ex:18-graph-and-cone}
设 \(k\) 为域。分别判断 \(k[x,y,z]_{(x,y,z)}/(z-xy)\) 与 \(k[x,y,z]_{(x,y,z)}/(z^2-xy)\) 是否正则，并写出它们的伴随分次环。
\end{exercise}
\begin{solution}
第一个环由消去 \(z\) 同构于 \(k[x,y]_{(x,y)}\)，维数和嵌入维数均为二，伴随分次环为 \(k[X,Y]\)。也可从最低项 \(z\) 得同一结果。
第二个环由三维局部整环除去非零因子，维数为二；方程在极大理想平方中，所以嵌入维数为三，不正则。由\cref{b3:prop:hypersurface-tangent-cone}，伴随分次环为 \(k[X,Y,Z]/(Z^2-XY)\)。这些判断对任意特征成立。
\end{solution}

\begin{exercise}\label{b3:ex:18-linear-parameter-change}
设 \(k\) 为域，\(R=k[x,y]_{(x,y)}\)。证明 \(x+y^2,y\) 是正则参数系，并比较商环 \(R/(x+y^2)\) 与 \(R/(x^2+y^3)\) 的正则性。
\end{exercise}
\begin{solution}
二者生成 \((x,y)\)，其一次像分别为 \(x,y\)，所以是极小生成组，因而正则参数系。第一个商同构于 \(k[y]_{(y)}\)，一维且正则。第二个商的非零方程位于平方理想中，所以维数为一而嵌入维数为二，不正则；即使某个特征下方程能进一步分解，这两个维数计算也不变。
\end{solution}

\begin{exercise}\label{b3:ex:18-mixed-characteristic-graded}
设 \(p\) 为素数，\(R=\mathbb Z[x]_{(p,x)}\)。证明 \(R\) 为二维正则局部环，并求其伴随分次环和剩余域的 Betti 数。
\end{exercise}
\begin{solution}
严格素链 \((0)\subsetneq(p)\subsetneq(p,x)\) 在局部化中保留，所以维数至少二。极大理想由 \(p,x\) 生成，高度定理使维数至多二。嵌入维数介于维数与生成元数之间，故也为二，环正则。由正文，伴随分次环为 \(\mathbb F_p[P,X]\)，剩余域 \(\mathbb F_p\) 的 Betti 数为 \((1,2,1)\)。此环的特征为零，剩余域特征却为 \(p\)，所以正则局部环不只来自某个域上的坐标原点。
\end{solution}

\begin{exercise}\label{b3:ex:18-parameter-powers-length}
设 \(k\) 为域，在 \(R=k[x,y]_{(x,y)}\) 中取正整数 \(a,b\)。证明 \(x^a,y^b\) 为正则序列，求商环的长度、投射维数和 Betti 数，并判断何时这个商环正则。
\end{exercise}
\begin{solution}
\(x^a\) 非零，故第一步单射；模 \(x^a\) 后，乘 \(y^b\) 仍单射，可按 \(1,x,\ldots,x^{a-1}\) 的 \(k[y]\) 系数逐项验证。商有基 \(x^iy^j\)、\(0\le i<a,0\le j<b\)，局部化不改变它，因此长度 \(ab\)。Koszul 消解极小，投射维数二，Betti 数 \((1,2,1)\)。商为零维局部环，只有 \(a=b=1\) 时是域并正则；否则 \(x\) 或 \(y\) 的某个非零像幂零。
\end{solution}

\begin{exercise}\label{b3:ex:18-regular-hypersurface-criterion}
设 \((R,\mathfrak m)\) 为 \(d>0\) 维交换 Noether 正则局部环，\(0\ne f\in\mathfrak m\)。证明 \(R/(f)\) 正则当且仅当 \(f\notin\mathfrak m^2\)。
\end{exercise}
\begin{solution}
\(R\) 是整环，故商的维数为 \(d-1\)。商的极大理想模平方为 \(\mathfrak m/(\mathfrak m^2+(f))\)。若 \(f\notin\mathfrak m^2\)，维数为 \(d-1\)；若 \(f\in\mathfrak m^2\)，维数仍为 \(d\)。与 Krull 维数比较即得等价。这个练习也给出上一题非正则性的另一种检验。
\end{solution}

\begin{exercise}\label{b3:ex:18-koszul-globaldim-three}
设 \(k\) 为域，令 \(R=k[x,y,z]_{(x,y,z)}\)、\(M=R/(x,y)\)。写出 \(M\) 的极小消解并求各 \(\operatorname{Tor}_i^R(M,k)\)。再比较 \(\operatorname{pd}_RM\) 与环的整体维数。
\end{exercise}
\begin{solution}
消解为 \(0\to R\xrightarrow{\binom{-y}{x}}R^2\xrightarrow{(x\quad y)}R\to M\to0\)。极小性使 Tor 各项依次为 \(k,k^2,k\)，更高次数为零。所以 \(\operatorname{pd}_RM=2\)，而 \(\operatorname{gldim}R=3\)。整体维数是所有模的上确界，不要求每个非自由模都达到这一上界。
\end{solution}

\begin{exercise}\label{b3:ex:18-abs-splitting-concrete}
设 \(k\) 为域，\(R=k[x,y]_{(x,y)}\)、\(\mathfrak m=(x,y)R\)、\(B=R/(x)\)。显式分解 \(\mathfrak m/x\mathfrak m\)，并写出到 \(\mathfrak m/xR\) 的截面。
\end{exercise}
\begin{solution}
\(\mathfrak m/x\mathfrak m\) 由 \(\bar x,\bar y\) 生成，其关系为 \(y\bar x=0\)；\(\bar y\) 所生成的部分自由秩一，而 \(\bar x\) 生成一个剩余域。因此
\(\mathfrak m/x\mathfrak m\cong k\oplus B\)。到 \(\mathfrak m/xR=(y)B\) 的映射杀掉第一项，并把第二项的基送到 \(y\)。截面将 \(by\) 送到 \(b\bar y\)。\(y\) 在 \(B\) 中非零因子保证表达唯一。
\end{solution}

\begin{exercise}\label{b3:ex:18-formal-regularity}
设 \(k\) 为域。证明 \(k[\![x,y,z]\!]\) 正则，并判断
\(k[\![x,y,z]\!]/(x+yz)\) 与 \(k[\![x,y,z]\!]/(x^2+yz)\) 是否正则。
\end{exercise}
\begin{solution}
幂级数环是坐标原点多项式局部环的完备化，由完备化正则性定理，它正则且维数三。两个方程均非零，所以商维数二。前者有非零一次项，商的嵌入维数二，故正则；后者全在平方理想中，商的嵌入维数三，故不正则。
\end{solution}

\begin{exercise}\label{b3:ex:18-localization-bases}
设 \(R\) 为交换 Noether 正则局部环，\(\mathfrak p\) 为高度一素理想。证明 \(R_{\mathfrak p}\) 为 DVR。解释为什么仅有这一结论尚未证明 \(\mathfrak p\) 在 \(R\) 中主生成。
\end{exercise}
\begin{solution}
局部化后正则，维数为 \(\operatorname{ht}\mathfrak p=1\)，又为整环，故为 DVR。但 DVR 的统一化元可以含不属于 \(\mathfrak p\) 的分母，只给出一个局部生成元；它也没有控制 \(\mathfrak p\) 在其他点的生成。Dedekind 整环的非主理想说明局部生成元不能无条件拼成全局生成元。本章额外用有限自由消解和行列式线完成了这一步。
\end{solution}

\begin{exercise}\label{b3:ex:18-stably-free-determinant}
设 \(A\) 为交换整环，\(r,n\ge0\) 为整数，\(P\) 为常秩 \(r\) 的有限生成投射 \(A\) 模，且 \(P\oplus A^n\cong A^{r+n}\)。证明 \(\det P\cong A\)，并说明 \(r=1\) 时得到什么。
\end{exercise}
\begin{solution}
直和最高外幂给 \(\det(P\oplus A^n)\cong\det P\otimes\det A^n\cong\det P\)，而右边自由模的最高外幂同构于 \(A\)。故 \(\det P\cong A\)。当 \(r=1\)，\(P=\det P\)，于是 \(P\) 自由。较高秩时这个推理只控制最高外幂，不能把 \(P\) 与 \(\det P\) 当成同一个模。
\end{solution}

\begin{exercise}\label{b3:ex:18-projective-not-ffr}
令 \(A=\mathbb Z[\sqrt{-5}]\)。利用\chapref{b3:ch:10}中的非主可逆理想 \(I=(2,1+\sqrt{-5})\subseteq A\)，说明“有限长度的投射消解”与“有限长度、各项有限秩的自由消解”在非局部环上不同。
\end{exercise}
\begin{solution}
\(I\) 自身是有限投射模，故投射维数零，已经具有长度零的投射消解。若另有有限自由消解，\cref{b3:lem:determinant-finite-free-resolution}会使秩一模 \(I\) 自由，即理想主生成，与前章已证的非主性矛盾。因此它没有有限自由消解。局部环上有限投射模自由，正是本章能够获得有限自由消解的原因。
\end{solution}

\begin{exercise}\label{b3:ex:18-higher-height-not-principal}
设 \(k\) 为域。证明二维正则局部环 \(R=k[x,y]_{(x,y)}\) 是唯一分解整环却不是主理想整环。给出不能用高度一判据使 \((x,y)\) 主生成的原因。
\end{exercise}
\begin{solution}
唯一分解性来自本章定理。若极大理想 \((x,y)\) 主生成，则其模平方的向量空间至多一维；实际上 \(x,y\) 的像构成二维基，矛盾。它的高度为二，故不在“高度一素理想主生成”的结论范围内。
\end{solution}

\begin{exercise}\label{b3:ex:18-regular-product}
取两个域 \(k_1,k_2\)。证明 \(A=k_1\times k_2\) 为正则环，且整体维数为零，但不是整环。
\end{exercise}
\begin{solution}
两个素理想的局部环分别为 \(k_1,k_2\)，所以都正则，\(A\) Noether，因而为正则环。对任意 \(A\) 模 \(M\)，幂等元 \(e_1=(1,0),e_2=(0,1)\) 给出 \(M=e_1M\oplus e_2M\)，两个分量分别为域上的向量空间。每个分量是对应投射模 \(Ae_i\) 的直和，故 \(M\) 投射，整体维数零。但 \(e_1e_2=0\) 且两者非零，所以不是整环。
\end{solution}

\begin{exercise}\label{b3:ex:18-inseparable-cotangent-sequence}
设 \(p\) 为素数，\(s\) 为 \(\mathbb F_p\) 上的不定元，令 \(k=\mathbb F_p(s)\)、\(A=k[T]\)、\(\mathfrak p=(T^p-s)\)、\(K=A/\mathfrak p\)。计算余切序列中
\(\mathfrak p/\mathfrak p^2\to\Omega_{A/k}\otimes_AK\to\Omega_{K/k}\to0\)
的各项和映射，并判断 \(A_{\mathfrak p}\) 的正则性。
\end{exercise}
\begin{solution}
\(\mathfrak p/\mathfrak p^2\) 以 \(T^p-s\) 的像为 \(K\) 基，故为一维；中项为 \(K\,\mathrm{d}T\)，末项为 \(K\,\mathrm{d}\bar T\)。第一箭头把生成元送到 \(\mathrm{d}(T^p-s)=0\)，为零；第二箭头为同构。另一方面 \(A_{\mathfrak p}\) 是一维主理想局部整环，极大理想由 \(T^p-s\) 生成，故正则。这说明中间微分模的纤维与极大理想模平方虽在此恰好同维，自然映射仍可能为零，不能照抄有理点的同构。
\end{solution}

\begin{exercise}\label{b3:ex:18-pure-inseparable-thickening}
设 \(p\) 为素数，\(s\) 为 \(\mathbb F_p\) 上的不定元，令 \(k=\mathbb F_p(s)\)、\(K=k[u]/(u^p-s)\)。求 \(K\otimes_kK\) 的长度、嵌入维数和深度，并判断其剩余域是否具有有限投射维数。
\end{exercise}
\begin{solution}
该环为 \(B=K[v]/(v^p)\)，其基为 \(1,v,\ldots,v^{p-1}\)，长度为 \(p\)。极大理想 \((v)\) 模平方维数一，故嵌入维数一；\(v^{p-1}\ne0\) 被整个极大理想零化，深度零。环维数零而嵌入维数一，不正则，ABS 判据使其剩余域投射维数无限。也可用交替乘 \(v\)、\(v^{p-1}\) 的周期极小消解验证。
\end{solution}

\begin{exercise}\label{b3:ex:18-smooth-base-matters}
设 \(k\) 为特征 \(p>0\) 的域。环 \(B=k[t]\) 作为 \(k\) 代数光滑；把它改看为 \(A=k[s]\) 代数，令 \(s\mapsto t^p\)。构造平方零提升障碍，证明后一映射不是形式光滑。
\end{exercise}
\begin{solution}
取 \(C=k[t,v]/(v^2)\)，让 \(A\) 通过 \(s\mapsto t^p+v\) 作用，令 \(J=(v)\)。模 \(J\) 后，恒等 \(B\to C/J=k[t]\) 与给定的 \(A\) 作用相容。若有提升，\(t\) 的像必为 \(t+v\cdot h(t)\)。其 \(p\) 次幂为 \(t^p\)，因为特征为 \(p\) 且 \(v^p=0\)。但 \(A\) 线性要求这个幂为 \(t^p+v\)，矛盾。底环改变了提升问题，虽然环 \(B\) 自身仍然正则。
\end{solution}
